\documentclass[instructions.tex]{subfiles}
\begin{document}
 
\chapter{De la correction fraternelle.}
 
\primo Corrigez-vous avant de corriger les autres; il sied mal de reprendre dans le prochain ce que l’on fait soi-même. Si vous avez du zèle contre les vices, que ce soit premièrement contre les vôtres. Réservez la miséricorde pour autrui, et la sévérité pour vous.
 
Si votre frère tombe dans une faute, dit Jésus-Christ, reprenez-le avec charité en secret, de peur de lui faire confusion devant les autres. S’il vous écoute, vous l’aurez gagné à Dieu; s’il ne vous écoute pas, s’il persévère dans le désordre, avertissez avec prudence ceux qui peuvent y apporter du remède. S’il pèche publiquement, reprenez-le publiquement, si sa faute est scandaleuse, et si vous jugez qu’une répréhension publique lui sera utile et aux autres. Personne n’étant obligé à une chose inutile, ne faites pas une correction, lorsque vous la prévoyez inutile, dangereuse et nuisible. Il faut souffrir avec patience ce que l’on ne peut corriger avec succès.\footnote{Detrahentem secretò proximo suo, hunc persequebar. Ps. 100. 5.}
 
\secundo N’entreprenez jamais de châtier et de punir les autres, que vous n’en ayez l’autorité. Père, mère, maître, supérieur, reprenez, menacez, châtiez: vous pêcheriez, si vous manquiez à ce devoir. On nuit aux bons lorsqu’on ne corrige pas les méchans \footnote{Bonis nocet qui malis parcit.}. On nuit aux méchans, même lorsqu’on ne les punit pas. La correction, dit le Sage, délivre un enfant de l’enfer. On nuit enfin à soi-même, on se rend complice du crime, quand on le laisse impuni. Il faut corriger, non pas tant parce qu’on a péché, que pour empêcher de pécher de nouveau, pour prévenir les désordres qui naissent de l’impunité.
 
La violence et la rigueur ne rendent pas toujours les corrections efficaces. Il y a des esprits que la vérité ne peut vaincre. Quand vous pileriez un imprudent dans un mortier, comme on pile le grain, vous ne lui ôteriez pas son imprudence et sa malice, dit le Sage\footnote{Prov. 27. 22.}. Il faut joindre la prière et la modération à la correction, recommander à Dieu ces esprits incorrigibles; Dieu seul est le médecin du cœur. Il faut avoir créé l’âme, dit saint Augustin, pour être capable de la guérir.
 
C’est de même perdre son temps que de prétendre corriger par la rigueur une femme querelleuse et farouche: C’est vouloir arrêter le vent, dit le Sage, et retenir dans sa main une poignée d’huile. Les maris doivent, en ce cas, profiter de l’avis que donne saint Paul: Chérissez vos femmes, et ne les traitez pas avec aigreur. On les gagne bien plutôt par une sage modération que par une indiscrète sévérité. Toute correction, pour être salutaire, doit être faite avec charité. Crier sans cesse dans sa famille, s’emporter de colère pour les moindres fautes, ce n’est pas corriger, mais irriter les esprits. Il est honteux, dit un ancien, de commettre une faute pour en punir une autre.
 
\tertio Craignez de perdre la confiance de vos inférieurs. Si, de peur d’être méprisé, vous voulez vous faire craindre, que ce soit par un zèle prudent, par la douceur, et non par la hauteur. Il vaut encore mieux, en quelque sorte, être méprisé que d’être haï. c’est un moindre mal d’affaiblir un peu son autorité par la douceur, que de la soutenir par trop de rigueur. Quiconque se fait obéir de la sorte, perd l’amitié de ses inférieurs. Il a peut-être autant d’ennemis à appréhender qu’il y a de gens qui le craignent.
 
\section{Résolutions}
 
\primo J’avertirai mon prochain, quand je croirai pouvoir le faire avec utilité; mais je choisirai le temps, l’occasion convenable, et je le ferai avec bonté et précaution, autant que les circonstances le permettront.
 
\secundo Je remplirai aussi le devoir de la correction envers mes inférieurs, mais toujours avec miséricorde, plutôt qu’avec dureté, et je n’userai de sévérité, qu’après avoir essayé vainement la douceur, ou quand je serai forcé d’en venir là pour réprimer un scandale public. 

\prayer{Mon Dieu, aidez-moi à remplir tous les devoirs que vous m’avez imposés envers mes semblables.}
 
\end{document}
